%%FILE blueprint/src/macros/common.tex
% In this file you should put all LaTeX macros and settings to be used both by
% the pdf version and the web version.
% This should be most of your macros.

% The theorem-like environments defined below are those that appear by default
% in the dependency graph. See the README of leanblueprint if you need help to
% customize this.
% The configuration below use the theorem counter for all those environments
% (this is what the [theorem] arguments mean) and never resets it.
% If you want for instance to number them within chapters then you can add
% [chapter] at the end of the next line.
\newtheorem{theorem}{Theorem}
\newtheorem{proposition}[theorem]{Proposition}
\newtheorem{lemma}[theorem]{Lemma}
\newtheorem{corollary}[theorem]{Corollary}

\newtheorem{metatheorem}[theorem]{Metatheorem}
\newtheorem{metametatheorem}[theorem]{Meta-metatheorem}

\newtheorem{example}{Example}
\newtheorem{remark}{Remark}

\theoremstyle{definition}
\newtheorem{definition}[theorem]{Definition}
\numberwithin{theorem}{chapter}

\newcommand{\N}{\mathbb{N}}
\newcommand{\R}{\mathbb{R}}
\newcommand{\Q}{\mathbb{Q}}
\newcommand{\Z}{\mathbb{Z}}
\newcommand{\formaleq}{\simeq}
\newcommand{\op}{\diamond}

\newcommand{\Equadef}[2]{%
  \expandafter\gdef\csname Equa#1\endcsname{\ensuremath{#2}}}
\newcommand{\Equareffull}[1]{\Equaref{#1} $(\csname Equa#1\endcsname)$}
\newcommand{\Equaref}[1]{\href{https://teorth.github.io/equational_theories/implications/?#1}{E#1}}

%%FILE blueprint/src/macros/web.tex
% In this file you should put macros to be used only by
% the web version. Of course they should have a corresponding
% version in macros/print.tex.
% Typically the printed version could have more fancy decorations.
% This will probably be a very short file.
%%FILE blueprint/src/macros/print.tex
% In this file you should put macros to be used only by
% the printed version. Of course they should have a corresponding
% version in macros/web.tex.
% Typically the printed version could have more fancy decorations.
% This should be a very short file.
%
% This file starts with dummy macros that ensure the pdf
% compiler will ignore macros provided by plasTeX that make
% sense only for the web version, such as dependency graph
% macros.


% Dummy macros that make sense only for web version.
\newcommand{\lean}[1]{}
\newcommand{\discussion}[1]{}
\newcommand{\leanok}{}
\newcommand{\mathlibok}{}
\newcommand{\notready}{}
% Make sure that arguments of \uses and \proves are real labels, by using invisible refs:
% latex prints a warning if the label is not defined, but nothing is shown in the pdf file.
% It uses LaTeX3 programming, this is why we use the expl3 package.
\ExplSyntaxOn
\NewDocumentCommand{\uses}{m}
 {\clist_map_inline:nn{#1}{\vphantom{\ref{##1}}}%
  \ignorespaces}
\NewDocumentCommand{\proves}{m}
 {\clist_map_inline:nn{#1}{\vphantom{\ref{##1}}}%
  \ignorespaces}
\ExplSyntaxOff

% https://tex.stackexchange.com/a/730164
% Workaround for broken interaction between cleveref and the 2025 LaTeX kernel
\AddToHook{env/proposition/begin}{\crefalias{theorem}{proposition}}
\AddToHook{env/lemma/begin}{\crefalias{theorem}{lemma}}
\AddToHook{env/corollary/begin}{\crefalias{theorem}{corollary}}
\AddToHook{env/metatheorem/begin}{\crefalias{theorem}{metatheorem}}
\AddToHook{env/metametatheorem/begin}{\crefalias{theorem}{metametatheorem}}
\AddToHook{env/example/begin}{\crefalias{theorem}{example}}
\AddToHook{env/remark/begin}{\crefalias{theorem}{remark}}
\AddToHook{env/definition/begin}{\crefalias{theorem}{definition}}
